%%%%%%%%%%%%%%%%%%%%%%%%%%%%%%%%%
%
%       EXERCÍCIO
%
%%%%%%%%%%%%%%%%%%%%%%%%%%%%%%%%%

\definevalorquestao

\begin{exercicioBanco}[\valorquestao]
Sejam \(A\) e \(B\) eventos de um espaço amostral. Responda cada uma das perguntas a seguir apresentando sua justificativa.

\begin{enumerate}[(a)]
    \item Se \(P(A) = \dfrac{1}{4}\) e \(P(B \mid A) = \dfrac{1}{2}\), então \(A\) e \(B\) são disjuntos?
    %
    \item Se \(P(A) = 0,3\), \(P(B) = 0,4\) e \(P(A\cap B) = 0,12\), então \(A\) e \(B\) são independentes?
\end{enumerate}

\end{exercicioBanco}


